\documentclass[12pt]{article}

% =========================
% 0) Metadata (EDIT ME)
% =========================
\newcommand{\CourseCode}{MATH 521}
\newcommand{\CourseTitle}{Analysis I}
\newcommand{\CourseSection}{LEC 002}
\newcommand{\Semester}{Spring 2026}
\newcommand{\Instructor}{Thierry Laurens}
\newcommand{\MeetingInfo}{MWF 1:20--2:10 \;|\; VV B119} % optional
\newcommand{\HomeworkNumber}{12} % <-- change each week
\newcommand{\YourName}{Alejandro De La Torre}
\newcommand{\DueDate}{April 24, 2026} % <-- or set e.g. January 31, 2026

% =========================
% 1) Core math + proof tools
% =========================
\usepackage{amsmath, amssymb, amsthm}

% =========================
% 2) Lists and spacing
% =========================
\usepackage{enumitem}
\setlist[enumerate,1]{label=\textbf{(\alph*)}, leftmargin=2em, itemsep=0.25em}
\setlist[itemize]{leftmargin=1.5em, itemsep=0.25em}
\setlength{\parskip}{0.35em}
\setlength{\parindent}{0pt}

% =========================
% 3) Page geometry + header/footer
% =========================
\usepackage{geometry}
\geometry{margin=1in}

\usepackage{fancyhdr}
\pagestyle{fancy}
\fancyhf{}
\lhead{\CourseCode\ --- \CourseSection, \Semester \;(\MeetingInfo)}
\rhead{Homework \HomeworkNumber}
\cfoot{\thepage}
\renewcommand{\headrulewidth}{0.4pt}
% Fixes common fancyhdr warning about header height:
\setlength{\headheight}{15pt}
% Optional: reclaim a bit of vertical space after increasing headheight
\addtolength{\topmargin}{-2pt}

% =========================
% 4) Links (subtle)
% =========================
\usepackage[hidelinks]{hyperref}
\setcounter{tocdepth}{1}

% =========================
% 5) Quotes / figures
% =========================
\usepackage{csquotes}
\usepackage{graphicx}
\graphicspath{{images/}} % put figures in ./images/

% =========================
% 6) Simple scratch box (no tcolorbox)
% =========================
% A lightweight environment for brief planning / scratch work.
% This avoids any tcolorbox dependencies.
\newenvironment{scratch}{%
  \par\smallskip
  \noindent\textbf{Discussion \& Scratch Work}\begin{quote}\small
}{%
  \end{quote}\smallskip
}

% =========================
% 7) Lightweight theorem-like environments (optional)
% =========================
\newtheorem*{claim}{Claim}
\newtheorem*{lemma}{Lemma}
\newtheorem*{remark}{Remark}
\newtheorem{theorem}{Theorem}
\newtheorem{proposition}{Proposition}
\newtheorem{corollary}{Corollary}
\newtheorem{definition}{Definition}
\newtheorem{example}{Example}
\newtheorem{exercise}{Exercise}

% =========================
% 8) Problem header command (with TOC + PDF bookmarks)
% Usage: \problem[Optional short title]{<number>}
% =========================
\makeatletter
\newcommand{\problem}[2][]{%
  \def\prob@title{Problem #2\if\relax\detokenize{#1}\relax\else\ (\textit{#1})\fi}%
  \phantomsection%
  \pdfbookmark[1]{\prob@title}{prob#2}%
  \addcontentsline{toc}{section}{\prob@title}%
  \section*{\prob@title}%
  \vspace{-0.25em}\hrule\vspace{0.8em}%
}
\makeatother

% =========================
% 9) Title block
% =========================
\title{Homework \HomeworkNumber}
\author{\YourName \\
\CourseCode\ --- \CourseTitle \\
\CourseSection \\
Instructor: \Instructor}
\date{\DueDate}

\begin{document}
\maketitle

% ------------------ collaboration + sources ------------------
% \section*{Collaboration Statement}
% Write “I worked alone” or list collaborators / external resources.
% "I worked on this homework independently. No outside collaborators or resources were used."

\tableofcontents
\bigskip
\hrule
\bigskip

\newpage

% ============================================================
% Problems (duplicate blocks as needed)
% ============================================================

\problem[]{1}
Let $(X,d_X)$ and $(Y,d_Y)$ be metric spaces, $f:X\to Y$, and $x_0\in X$. Prove that
$f$ is continuous at $x_0$ if and only if for any open set $V\subseteq Y$ containing
$f(x_0)$, there is an open set $U\subseteq X$ containing $x_0$ so that $x\in U$ implies
$f(x)\in V$.

\begin{scratch}
This problem is asking me to prove that two definitions of continuity at a point are equivalent. From lecture, I know the standard $\varepsilon$--$\delta$ definition: $f$ is continuous at $x$ if for every $\varepsilon>0$ there exists $\delta>0$ such that
\[
d_X(x',x)<\delta \implies d_Y(f(x'),f(x))<\varepsilon.
\]
I also know the open-set formulation: $f$ is continuous at $x$ if for every open set $V\subseteq Y$ with $f(x)\in V$, there exists an open set $U\subseteq X$ such that
\[
x\in U
\qquad\text{and}\qquad
f(U)\subseteq V.
\]

So the goal is to show these two definitions are equivalent. Since this is an ``if and only if,'' I prove both directions.

For the forward direction, I assume the $\varepsilon$--$\delta$ definition. If $V$ is open and contains $f(x)$, then there exists an open ball $B_\varepsilon(f(x))\subseteq V$. Then continuity gives a $\delta>0$ such that $f(B_\delta(x))\subseteq B_\varepsilon(f(x))\subseteq V$. So I take $U=B_\delta(x)$.

For the reverse direction, I assume the open-set condition. Given $\varepsilon>0$, I consider the open ball $B_\varepsilon(f(x))$. By hypothesis, there exists an open set $U$ containing $x$ such that $f(U)\subseteq B_\varepsilon(f(x))$. Since $U$ is open, there exists $\delta>0$ such that $B_\delta(x)\subseteq U$. This gives the $\varepsilon$--$\delta$ condition.
\end{scratch}

\textbf{Proof.}\\
$(\Rightarrow)$ Assume that $f$ is continuous at $x\in X$. Let $V\subseteq Y$ be open with $f(x)\in V$. Then there exists $\varepsilon>0$ such that
\[
B_\varepsilon(f(x))\subseteq V.
\]
By continuity at $x$, there exists $\delta>0$ such that
\[
d_X(x',x)<\delta \implies d_Y(f(x'),f(x))<\varepsilon.
\]
Let $U:=B_\delta(x)$. Then $U$ is open, $x\in U$, and if $x'\in U$, then $f(x')\in B_\varepsilon(f(x))\subseteq V$. Hence
\[
f(U)\subseteq V.
\]

$(\Leftarrow)$ Assume that for every open set $V\subseteq Y$ with $f(x)\in V$, there exists an open set $U\subseteq X$ such that $x\in U$ and $f(U)\subseteq V$.

Let $\varepsilon>0$ and set
\[
V:=B_\varepsilon(f(x)).
\]
Then $V$ is open and contains $f(x)$, so there exists an open set $U\subseteq X$ such that
\[
x\in U
\qquad\text{and}\qquad
f(U)\subseteq B_\varepsilon(f(x)).
\]
Since $U$ is open and $x\in U$, there exists $\delta>0$ such that
\[
B_\delta(x)\subseteq U.
\]
Now if $d_X(x',x)<\delta$, then $x'\in U$, so
\[
f(x')\in B_\varepsilon(f(x)),
\]
which implies
\[
d_Y(f(x'),f(x))<\varepsilon.
\]
Therefore $f$ is continuous at $x$.
\qed

\newpage

\problem[]{2}
Let $(X,d_X)$ and $(Y,d_Y)$ be metric spaces and $f:X\to Y$ a function. Prove
that $f$ is continuous if and only if $f(\overline{A})\subseteq \overline{f(A)}$ for any
$A\subseteq X$.

\begin{scratch}
This problem is asking me to prove an equivalent characterization of continuity using closures. I am given a function $f:X\to Y$ between metric spaces, and I need to show that
\[
f \text{ is continuous } \iff f(\overline{A}) \subseteq \overline{f(A)} \quad \text{for all } A\subseteq X.
\]

So I need to prove both directions.

First, I unpack what this statement means. The set $\overline{A}$ consists of all points $x$ that are either in $A$ or are limits of sequences from $A$. Similarly, $\overline{f(A)}$ consists of all points in $Y$ that are limits of sequences from $f(A)$.

So the statement $f(\overline{A}) \subseteq \overline{f(A)}$ is saying:
if $x$ is either in $A$ or is a limit of points in $A$, then $f(x)$ should be either in $f(A)$ or a limit of points in $f(A)$.

This strongly suggests using the sequential characterization of closure and the fact from lecture that:
\[
f \text{ continuous } \iff x_n \to x \Rightarrow f(x_n)\to f(x).
\]

For the forward direction, I assume $f$ is continuous. Then if $x\in \overline{A}$, I can find a sequence $a_n\in A$ with $a_n\to x$. Continuity will give $f(a_n)\to f(x)$, which means $f(x)\in \overline{f(A)}$.

For the reverse direction, I should not use subsequences. Instead, I will use the open-set characterization from Problem 1. So I fix an open set $V\subseteq Y$ and a point $x\in X$ such that $f(x)\in V$, and I define
\[
A:=X\setminus f^{-1}(V).
\]
If $x$ were in $\overline{A}$, then the hypothesis would imply
\[
f(x)\in f(\overline{A})\subseteq \overline{f(A)}.
\]
But every point of $f(A)$ lies in $Y\setminus V$, so $\overline{f(A)}\subseteq Y\setminus V$ because $Y\setminus V$ is closed. That would force $f(x)\notin V$, a contradiction. Hence $x\notin \overline{A}$.

Since $x\notin \overline{A}$, there exists an open set $U$ containing $x$ with $U\cap A=\varnothing$. Therefore $U\subseteq X\setminus A=f^{-1}(V)$, so $f(U)\subseteq V$. This is exactly the open-set characterization of continuity.
\end{scratch}

\textbf{Proof.}\\
$(\Rightarrow)$ Assume that $f$ is continuous. Let $A\subseteq X$, and let $x\in \overline{A}$. By the sequential characterization of closure, there exists a sequence $\{a_n\}\subseteq A$ such that
\[
a_n \to x.
\]
Since $f$ is continuous, we have
\[
f(a_n)\to f(x).
\]
Each $f(a_n)\in f(A)$, so $f(x)$ is a limit of points in $f(A)$. Hence
\[
f(x)\in \overline{f(A)}.
\]
Therefore
\[
f(\overline{A}) \subseteq \overline{f(A)}.
\]

$(\Leftarrow)$ Assume that for all $A\subseteq X$, we have
\[
f(\overline{A}) \subseteq \overline{f(A)}.
\]
We show that $f$ is continuous.

Let $V\subseteq Y$ be open, and let $x\in X$ satisfy
\[
f(x)\in V.
\]
Set
\[
A:=X\setminus f^{-1}(V).
\]
We claim that $x\notin \overline{A}$. Suppose for contradiction that
\[
x\in \overline{A}.
\]
Then, since $x\in \overline{A}$, we have $f(x)\in f(\overline{A})$. By the hypothesis,
\[
f(\overline{A})\subseteq \overline{f(A)},
\]
so
\[
f(x)\in \overline{f(A)}.
\]

Now if $y\in f(A)$, then $y=f(a)$ for some $a\in A$. Since $a\notin f^{-1}(V)$, we must have
\[
f(a)\notin V.
\]
Thus
\[
f(A)\subseteq Y\setminus V.
\]
Because $V$ is open, the set $Y\setminus V$ is closed. Therefore
\[
\overline{f(A)}\subseteq Y\setminus V.
\]
It follows that
\[
f(x)\in Y\setminus V,
\]
which contradicts $f(x)\in V$.

Hence $x\notin \overline{A}$. Since $\overline{A}$ is closed, its complement is open, so there exists an open set $U\subseteq X$ such that
\[
x\in U
\qquad\text{and}\qquad
U\cap A=\varnothing.
\]
Therefore
\[
U\subseteq X\setminus A=f^{-1}(V).
\]
So if $u\in U$, then $u\in f^{-1}(V)$, which means $f(u)\in V$. Hence
\[
f(U)\subseteq V.
\]

We have shown that for every open set $V\subseteq Y$ containing $f(x)$, there exists an open set $U\subseteq X$ containing $x$ such that $f(U)\subseteq V$. By Problem 1, this means that $f$ is continuous at $x$. Since $x\in X$ was arbitrary, $f$ is continuous on $X$.
\qed

\newpage

\problem[]{3}
Let $f:\mathbb{R}\to\mathbb{R}$ be the function $f(x)=\frac{1}{1+x^2}$.
\begin{enumerate}[label=\textbf{(\alph*)}]
\item Prove that $f$ is continuous.
\item Find an open set $U\subseteq \mathbb{R}$ so that $f(U)$ is not open.
\item Find a closed set $F\subseteq \mathbb{R}$ so that $f(F)$ is not closed.
\item Find a set $A\subseteq \mathbb{R}$ so that $f(\overline{A})\neq \overline{f(A)}$.
\end{enumerate}

\begin{scratch}
This problem is asking me to analyze a specific function
\[
f(x)=\frac{1}{1+x^2}
\]
from several perspectives: continuity, behavior of images of open/closed sets, and how closure interacts with the function.

For part (a), I need to prove $f$ is continuous. From Lecture 30, I know I can use the sequential definition: $f$ is continuous if $x_n\to x$ implies $f(x_n)\to f(x)$. This is the easiest approach here since rational functions behave well under limits.

For part (b), I need to find an open set $U$ such that $f(U)$ is not open. So I need to recall that a set is not open if it fails to contain an open ball around some point. Intuitively, $f$ maps $\mathbb{R}$ into $(0,1]$, and the maximum value $1$ is only achieved at $x=0$. That suggests that images of open sets containing $0$ will include $1$ but not values larger than $1$, which creates a boundary point—so the image won’t be open.

For part (c), I need a closed set $F$ such that $f(F)$ is not closed. From Lecture 31, continuous functions do not necessarily preserve closed sets unless the domain is compact. So I should look for a closed set whose image “misses” a limit point.

For part (d), I need to find a set $A$ where
\[
f(\overline{A}) \neq \overline{f(A)}.
\]
From Problem 2, I know that continuity gives only one inclusion:
\[
f(\overline{A}) \subseteq \overline{f(A)}.
\]
So I need an example where this inclusion is strict. This typically happens when limit points are “collapsed” under the function, meaning some points in $\overline{f(A)}$ are not hit by $f(\overline{A})$.
\end{scratch}

\textbf{Proof.}

\textbf{(a)} Let $\{x_n\}$ be a sequence in $\mathbb{R}$ such that $x_n \to x$. Then
We use the standard limit laws from earlier lectures: polynomials and reciprocal functions are continuous where defined. Then
\[
x_n^2 \to x^2,
\]
so
\[
1+x_n^2 \to 1+x^2.
\]
Since the function $t \mapsto \frac{1}{t}$ is continuous for $t\neq 0$, we have
\[
\frac{1}{1+x_n^2} \to \frac{1}{1+x^2}.
\]
Thus $f(x_n)\to f(x)$, so $f$ is continuous.

\textbf{(b)} Let $U=\mathbb{R}$, which is open. Then
\[
f(U)=\left(0,1\right].
\]
This set is not open in $\mathbb{R}$ because $1\in f(U)$, but there is no open interval around $1$ contained in $(0,1]$. Hence $f(U)$ is not open.

\textbf{(c)} Let $F=[0,\infty)$, which is closed in $\mathbb{R}$. Then
\[
f(F)=\left(0,1\right].
\]
As above, $(0,1]$ is not closed because it does not contain $0$, which is a limit point. Hence $f(F)$ is not closed.

\textbf{(d)} Let
\[
A=(0,\infty).
\]
Then
\[
\overline{A}=[0,\infty).
\]
So
\[
f(\overline{A})=f([0,\infty))=(0,1].
\]
On the other hand,
\[
f(A)=f((0,\infty))=(0,1),
\]
so
\[
\overline{f(A)}=[0,1].
\]
Thus
\[
f(\overline{A})=(0,1] \neq [0,1]=\overline{f(A)}.
\]
\qed

\newpage

\problem[]{4}
Consider the metric spaces $(\mathbb{R},|\cdot|)$ and $(\mathbb{R}^2,d_2)$.
\begin{enumerate}[label=\textbf{(\alph*)}]
\item Prove that the functions $f,g:\mathbb{R}^2\to\mathbb{R}$ given by $f(x,y)=x$ and
$g(x,y)=y$ are continuous. (Hint: Use sequences.)
\item Prove that if $h:[0,1]\to\mathbb{R}^2$ is a continuous function such that
$h(0)=(x_0,y_0)$, $h(1)=(x_1,y_1)$, and $x_0<0<x_1$, then there exists a point
$t\in(0,1)$ so that $h(t)$ lies on the $y$-axis.
\end{enumerate}

\begin{scratch}
For part (a), the problem explicitly tells me to use sequences, so I should use the sequential characterization of continuity from lecture: a function is continuous if whenever a sequence converges in the domain, its image converges in the codomain. Here the functions are just the coordinate projections
\[
f(x,y)=x
\qquad\text{and}\qquad
g(x,y)=y.
\]
So if $(x_n,y_n)\to (x,y)$ in $(\mathbb{R}^2,d_2)$, then by the definition of the Euclidean metric,
\[
d_2((x_n,y_n),(x,y))=\sqrt{(x_n-x)^2+(y_n-y)^2}\to 0.
\]
That immediately forces both $x_n\to x$ and $y_n\to y$, so the coordinate functions should be continuous.

For part (b), the key idea is to compose the given path $h:[0,1]\to \mathbb{R}^2$ with the projection onto the first coordinate. If I define
\[
\varphi(t):=f(h(t)),
\]
then $\varphi:[0,1]\to \mathbb{R}$ is continuous because $h$ is continuous and part (a) says $f$ is continuous. Also,
\[
\varphi(0)=x_0<0
\qquad\text{and}\qquad
\varphi(1)=x_1>0.
\]
So now I can use the Intermediate Value Theorem from calculus on $\varphi$. That gives some $t\in(0,1)$ with $\varphi(t)=0$, which means the first coordinate of $h(t)$ is $0$. In other words, $h(t)$ lies on the $y$-axis.
\end{scratch}

\textbf{Proof.}

\textbf{(a)} We first prove that $f:\mathbb{R}^2\to\mathbb{R}$ given by $f(x,y)=x$ is continuous. Let $\{(x_n,y_n)\}$ be a sequence in $\mathbb{R}^2$ such that
\[
(x_n,y_n)\to (x,y)
\]
in $(\mathbb{R}^2,d_2)$. Then
\[
d_2((x_n,y_n),(x,y))=\sqrt{(x_n-x)^2+(y_n-y)^2}\to 0.
\]
Since
\[
|x_n-x|\le \sqrt{(x_n-x)^2+(y_n-y)^2}=d_2((x_n,y_n),(x,y)),
\]
it follows by the squeeze theorem that
\[
x_n\to x.
\]
Hence
\[
f(x_n,y_n)=x_n\to x=f(x,y).
\]
Therefore $f$ is continuous.

The proof for $g:\mathbb{R}^2\to\mathbb{R}$ given by $g(x,y)=y$ is exactly the same. If $(x_n,y_n)\to (x,y)$, then
\[
|y_n-y|\le \sqrt{(x_n-x)^2+(y_n-y)^2}=d_2((x_n,y_n),(x,y))\to 0,
\]
so
\[
y_n\to y.
\]
Thus
\[
g(x_n,y_n)=y_n\to y=g(x,y),
\]
and therefore $g$ is continuous.

\textbf{(b)} Let $h:[0,1]\to\mathbb{R}^2$ be continuous, with
\[
h(0)=(x_0,y_0),
\qquad
h(1)=(x_1,y_1),
\qquad
x_0<0<x_1.
\]
Define
\[
\varphi:[0,1]\to\mathbb{R},
\qquad
\varphi(t):=f(h(t)),
\]
where $f:\mathbb{R}^2\to\mathbb{R}$ is the first-coordinate projection from part (a), namely $f(x,y)=x$.

Since $h$ is continuous and $f$ is continuous, the composition $\varphi=f\circ h$ is continuous on $[0,1]$. Moreover,
\[
\varphi(0)=f(h(0))=f(x_0,y_0)=x_0<0,
\]
and
\[
\varphi(1)=f(h(1))=f(x_1,y_1)=x_1>0.
\]
By the Intermediate Value Theorem from calculus, which we may assume, there exists some $t\in(0,1)$ such that
\[
\varphi(t)=0.
\]
But $\varphi(t)=f(h(t))$ is exactly the first coordinate of $h(t)$. So the first coordinate of $h(t)$ is $0$, which means
\[
h(t)\in \{(0,y):y\in\mathbb{R}\},
\]
that is, $h(t)$ lies on the $y$-axis.
\qed

\newpage

\problem[]{5}
Going forward, you may assume (no proof necessary) any fact about $\sin x$ and
$\cos x$ from calculus, such as $\sin,\cos:\mathbb{R}\to[-1,1]$ are continuous functions.
Prove that the function $f:\mathbb{R}^2\to\mathbb{R}$ given by
\[
f(x,y)=x^2+y^2+\sin(x-y)
\]
has a global minimum. (Hint: Notice that we cannot directly apply the extreme value
theorem here. However, we can apply it to any compact subset of $\mathbb{R}^2$.)

\begin{scratch}
The goal is to prove that
\[
f(x,y)=x^2+y^2+\sin(x-y)
\]
has a global minimum on $\mathbb{R}^2$. The hint is important: I cannot apply the Extreme Value Theorem directly on all of $\mathbb{R}^2$, because $\mathbb{R}^2$ is not compact. So the idea is to first show that I only need to look on some compact subset of $\mathbb{R}^2$, and then apply the Extreme Value Theorem there.

First, I note that $f$ is continuous. The functions $(x,y)\mapsto x^2+y^2$ and $(x,y)\mapsto x-y$ are continuous, and by the problem statement I may assume that $\sin$ is continuous. Therefore $(x,y)\mapsto \sin(x-y)$ is continuous, and so $f$ is continuous as a sum of continuous functions.

Next, I want a lower bound that tells me what happens far away from the origin. Since
\[
-1\le \sin(x-y)\le 1,
\]
I get
\[
f(x,y)\ge x^2+y^2-1.
\]
So if $(x,y)$ is far from the origin, then $x^2+y^2$ is large, and therefore $f(x,y)$ is large as well. This is the key coercive-type estimate.

To use this effectively, I want to compare that lower bound with the value of $f$ at some specific point. A convenient choice is
\[
\left(-\frac12,\frac12\right),
\]
because then
\[
x-y=-1,
\]
so
\[
f\!\left(-\frac12,\frac12\right)
= \frac14+\frac14+\sin(-1)
= \frac12-\sin(1)<0.
\]
On the other hand, if $x^2+y^2\ge 1$, then
\[
f(x,y)\ge x^2+y^2-1\ge 0.
\]
So every point outside the closed unit disk has $f(x,y)\ge 0$, while the point $\left(-\frac12,\frac12\right)$ inside the closed unit disk gives a negative value. That means a global minimizer cannot occur outside the closed unit disk.

So I restrict to
\[
K=\{(x,y)\in\mathbb{R}^2:x^2+y^2\le 1\}.
\]
This set is closed and bounded in $\mathbb{R}^2$, hence compact. Since $f$ is continuous, the Extreme Value Theorem implies that $f$ attains a minimum on $K$. Because every point outside $K$ has value at least $0$, while some point in $K$ has value $<0$, that minimum on $K$ is automatically a global minimum on all of $\mathbb{R}^2$.
\end{scratch}

\textbf{Proof.}\\
Define
\[
f:\mathbb{R}^2\to\mathbb{R},
\qquad
f(x,y)=x^2+y^2+\sin(x-y).
\]

First, $f$ is continuous on $\mathbb{R}^2$. Indeed, $(x,y)\mapsto x^2+y^2$ is continuous, $(x,y)\mapsto x-y$ is continuous, and $\sin:\mathbb{R}\to\mathbb{R}$ is continuous, so $(x,y)\mapsto \sin(x-y)$ is continuous. Therefore $f$ is continuous as a sum of continuous functions.

Next, since
\[
\sin(x-y)\ge -1
\]
for all $(x,y)\in\mathbb{R}^2$, we have
\[
f(x,y)=x^2+y^2+\sin(x-y)\ge x^2+y^2-1.
\]
In particular, whenever $x^2+y^2\ge 1$, we get
\[
f(x,y)\ge 0.
\]

Now consider the point
\[
p=\left(-\frac12,\frac12\right).
\]
Then
\[
f(p)
=
f\!\left(-\frac12,\frac12\right)
=
\frac14+\frac14+\sin\!\left(-1\right)
=
\frac12-\sin(1).
\]
Since $\sin(1)>\frac12$, it follows that
\[
f(p)<0.
\]
Also, $p$ lies in the closed unit disk
\[
K:=\{(x,y)\in\mathbb{R}^2:x^2+y^2\le 1\}.
\]

The set $K$ is closed and bounded in $\mathbb{R}^2$, hence compact. Since $f$ is continuous, the Extreme Value Theorem implies that there exists some $q\in K$ such that
\[
f(q)\le f(z)\qquad \text{for all } z\in K.
\]

We claim that $q$ is a global minimizer of $f$ on $\mathbb{R}^2$. Let $z=(x,y)\in\mathbb{R}^2$.

If $z\in K$, then by definition of $q$,
\[
f(q)\le f(z).
\]

If $z\notin K$, then $x^2+y^2>1$, so from the estimate above,
\[
f(z)\ge 0.
\]
But
\[
f(q)\le f(p)<0,
\]
so again
\[
f(q)<0\le f(z).
\]

Thus in every case,
\[
f(q)\le f(z)\qquad \text{for all } z\in\mathbb{R}^2.
\]
Therefore $f$ has a global minimum on $\mathbb{R}^2$.
\qed

\newpage

\problem[]{6}
Let $X$ and $Y$ be metric spaces and $f:X\to Y$ a function. Prove that if $X$ is
compact and $f$ is continuous and bijective, then the inverse function $f^{-1}:Y\to X$ is also
continuous. (Hint: It suffices to show that the preimage of any closed subset of $X$ under $f^{-1}$
is closed in $Y$.)

\begin{scratch}
The goal is to prove that $f^{-1}:Y\to X$ is continuous. The hint is telling me exactly which continuity criterion to use: it is enough to show that for every closed set $C\subseteq X$, the preimage of $C$ under $f^{-1}$ is closed in $Y$.

So I should start with an arbitrary closed set $C\subseteq X$ and understand what
\[
(f^{-1})^{-1}(C)
\]
actually is. Since $f$ is bijective, this set is just
\[
f(C).
\]
So the real task is to prove that $f(C)$ is closed in $Y$ whenever $C$ is closed in $X$.

Now the compactness assumption is what makes this work. Since $X$ is compact and $C$ is closed in $X$, I know from lecture that $C$ is compact. Since $f$ is continuous, the image of a compact set under $f$ is compact. And since $Y$ is a metric space, compact subsets of $Y$ are closed. So the chain is:
\[
C \text{ closed in } X
\;\Longrightarrow\;
C \text{ compact}
\;\Longrightarrow\;
f(C) \text{ compact}
\;\Longrightarrow\;
f(C) \text{ closed in } Y.
\]
That is exactly what I need to show that $f^{-1}$ is continuous.
\end{scratch}

\textbf{Proof.}\\
We prove that $f^{-1}:Y\to X$ is continuous by using the closed-set criterion.

Let $C\subseteq X$ be closed. Since $X$ is compact and $C$ is closed in $X$, it follows that $C$ is compact.

Because $f:X\to Y$ is continuous, the image of a compact set under $f$ is compact. Hence
\[
f(C)
\]
is compact in $Y$.

Now $Y$ is a metric space, and in a metric space every compact set is closed. Therefore $f(C)$ is closed in $Y$.

Since $f$ is bijective, we have
\[
(f^{-1})^{-1}(C)=f(C).
\]
Thus $(f^{-1})^{-1}(C)$ is closed in $Y$ for every closed set $C\subseteq X$.

Therefore, by the closed-set criterion for continuity, $f^{-1}$ is continuous.
\qed

% ============================================================
% Optional appendix: keep a personal toolbox of definitions/theorems
% ============================================================
\newpage
\appendix
\section*{Appendix: Toolbox}
\addcontentsline{toc}{section}{Appendix: Toolbox}

% Example (using amsthm environments instead of boxed tcolorboxes):
% \begin{definition}
% A metric space is ...
% \end{definition}
%
% \begin{theorem}[Heine--Borel]
% ...
% \end{theorem}

\end{document}
